\chapter{Relaciones}

\section{Definiciones básicas}

\begin{defn}
Una relación binaria en un tipo \lean{X} es una función \lean{X → X → Prop},
es decir, una proposición sobre pares de elementos de \lean{X}.
\begin{leancode}
def relation (X : Type) := X → X → Prop
\end{leancode}
\end{defn}

\begin{defn}
Sea \lean{r} una relación en \lean{A}. Entonces
\lean{r} induce las siguientes relaciones binarias.

\begin{enumerate}[(i)]
\item \lean{ReflGen r} es la cerradura reflexiva de \lean{r}, 
    definida inductivamente de la siguiente manera:
\begin{leancode}
inductive ReflGen {A : Type} (r : relation A) : relation A
  | refl {a : A}: ReflGen r a a
  | single {a b : A} : r a b → ReflGen r a b
\end{leancode}
\item \lean{ReflTransGen r} es la cerradura reflexiva y transitiva de \lean{r},
    definida inductivamente de la siguiente manera:
\begin{leancode}
inductive ReflTransGen {A : Type} (r : relation A) : relation A
  | refl {a : A} : ReflTransGen r a a
  | tail {a b c : A} : ReflTransGen r a b → r b c → ReflTransGen r a c
\end{leancode}
\end{enumerate}
\end{defn}

De ahora en adelante escribimos \lean{variable {A : Type}},
para evitar la repetición de \lean{{A : Type}}. Mientras que
 en los diagramas escribimos $\to$ para \lean{r a b},
$\to^=$ para \lean{ReflGen r a b} y
$\twoheadrightarrow$ para \lean{ReflTransGen r a b}.

\begin{defn}
Sea \lean{r} una relación en \lean{A}.
\begin{enumerate}[(i)]
\item \lean{Join r} es una nueva relación en la que 
dos términos están relacionados 
si existe un tercer término con el que ambos se relacionen.

\begin{leancode}
def Join (r : relation A) : relation A := fun x y ↦ ∃ z, r x z ∧ r y z
\end{leancode}

\item Una relación tiene la \textit{propiedad del diamante}
cuando todas las reducciones con un origen común se pueden unir.

\begin{center}
\begin{tikzpicture}
  \node (a) at (0, 2) {\lean{a}};
  \node (b) at (2, 2) {\lean{b}};
  \node (c) at (0, 0) {\lean{c}};
  \node (d) at (2, 0) {\lean{d}};
  \draw[->]           (a) -- node[above] {\lean{r}} (b);
  \draw[->]           (a) -- node[left]  {\lean{r}} (c);
  \draw[->, dashed] (b) -- node[right] {\lean{r}} (d);
  \draw[->, dashed] (c) -- node[below] {\lean{r}} (d);
\end{tikzpicture}
\end{center}

\begin{leancode}
def Diamond (r : relation A) := ∀ {a b c : A}, r a b → r a c → Join r b c
\end{leancode}
\end{enumerate}
\end{defn}

\subsection{Confluencia y Conmutación}

\begin{defn}
 Una relación es \textbf{confluente} cuando su cerradura 
 reflexiva y transitiva tiene la propiedad del diamante.

 \begin{center}
\begin{tikzpicture}
  \node (a) at (0, 2) {\lean{a}};
  \node (b) at (2, 2) {\lean{b}};
  \node (c) at (0, 0) {\lean{c}};
  \node (d) at (2, 0) {\lean{d}};
  \draw[->>]           (a) -- node[above] {\lean{r}} (b);
  \draw[->>]           (a) -- node[left]  {\lean{r}} (c);
  \draw[->>, dashed] (b) -- node[right] {\lean{r}} (d);
  \draw[->>, dashed] (c) -- node[below] {\lean{r}} (d);
\end{tikzpicture}
\end{center}

\begin{leancode}
def Confluent (r : relation A) := Diamond (ReflTransGen r)
\end{leancode}
\end{defn}

\begin{thm}
Si una relación tiene la propiedad del diamante, entonces es confluente.
\begin{leancode}
theorem Diamond.toConfluent (h : Diamond r) : Confluent r 
\end{leancode}
\end{thm}

\begin{defn}
Sean \lean{r₁} y \lean{r₂} dos relaciones en \lean{A}.

\begin{enumerate}[(i)]

\item \lean{r₁} y \lean{r₂} \textbf{conmutan} si:

\begin{center}
\begin{tikzpicture}
  \node (A) at (0, 2) {\lean{x}};
  \node (B) at (2, 2) {\lean{y₁}};
  \node (C) at (0, 0) {\lean{y₂}};
  \node (D) at (2, 0) {\lean{z}};
  \draw[->>]         (A) -- node[above] {\lean{r₁}} (B);
  \draw[->>]         (A) -- node[left]  {\lean{r₂}} (C);
  \draw[->>, dashed] (B) -- node[right] {\lean{r₂}} (D);
  \draw[->>, dashed] (C) -- node[below] {\lean{r₁}} (D);
\end{tikzpicture}
\end{center}

\begin{leancode}
def Commute (r₁ r₂ : relation A) := 
  ∀ {x y₁ y₂}, ReflTransGen r₁ x y₁ → ReflTransGen r₂ x y₂ → 
    ∃ z, ReflTransGen r₂ y₁ z ∧ ReflTransGen r₁ y₂ z
\end{leancode}


\item \lean{r₁} y \lean{r₂} \textbf{conmutan fuertemente} si:

\begin{center}
\begin{tikzpicture}
  \node (A) at (0, 2) {\lean{x}};
  \node (B) at (2, 2) {\lean{y₁}};
  \node (C) at (0, 0) {\lean{y₂}};
  \node (D) at (2, 0) {\lean{z}};
  \draw[->] (A) -- node[above] {\lean{r₁}}  (B);
  \draw[->] (A) -- node[left] {\lean{r₂}} (C);
  \draw[->, dashed] (B)-- node[right] {\lean{r₂}} node[left] {$=$} (D);
  \draw[->>, dashed] (C)-- node[below]  {\lean{r₁}}  (D);
\end{tikzpicture}
\end{center}

\begin{leancode}
def StronglyCommute (r₁ r₂ : relation A) :=
  ∀ {x y₁ y₂}, r₁ x y₁ → r₂ x y₂ → ∃ z, ReflGen r₂ y₁ z ∧ ReflTransGen r₁ y₂ z
\end{leancode}

\end{enumerate}

\end{defn}

\begin{thm}
Si \lean{r₁} y \lean{r₂} conmutan fuertemente, entonces también conmutan.
\begin{leancode}
theorem StronglyCommute.toCommute {r₁ r₂ : relation A} (h : StronglyCommute r₁ r₂) : 
  Commute r₁ r₂ 
\end{leancode}
\end{thm}

\begin{thm}
La conmutación entre relaciones es simétrica, es decir, 
si \lean{r₁} y \lean{r₂} conmutan, 
entonces \lean{r₂} y \lean{r₁} también.
\begin{leancode}
theorem Commute.symmetric : Symmetric (@Commute A)
\end{leancode}
\end{thm}


